%=============================================================================!
% 16.9  FREQUENCIES IN A SIGNAL                                               !
%=============================================================================!
\section{Frequencies in a signal}
\label{sec:ob-fourier}

A chord struck on a piano reaches the ear as one wobbling pressure, yet a
listener hears two notes, and a tuner's display shows two peaks at their
pitches. Something takes a signal in time apart into the frequencies it
holds. The same operation turns the velocities of a run into the
frequencies its atoms vibrate at, and it has been used twice already
without explanation, to compute correlation functions quickly in
Chapter~15 and in \cref{sec:ob-vacf}. This section supplies it: complex
numbers, the Fourier transform, its discrete form, the fast algorithm,
and what sampling a signal at intervals does to the frequencies one can
see.

\begin{toolbox}[Complex numbers and Euler's formula]
A complex number $z = x + \mathrm{i}y$ is a pair of real numbers $x$ and
$y$ combined with a symbol $\mathrm{i}$ whose square is $-1$. They add and
multiply as ordinary numbers do, replacing $\mathrm{i}^2$ by $-1$
wherever it appears: $(a + \mathrm{i}b)(c + \mathrm{i}d) = ac - bd +
\mathrm{i}(ad + bc)$. Drawn as the point $(x, y)$, $z$ is an arrow from
the origin; its length is $\lvert z\rvert = \sqrt{x^2 + y^2}$, and with
the conjugate $\bar z = x - \mathrm{i}y$, $z\bar z = x^2 + y^2 = \lvert z
\rvert^2$. The Taylor series of \cref{sec:mo-taylor} for $\mathrm{e}^x$,
$\cos x$ and $\sin x$ are
\begin{equation*}
   \mathrm{e}^x = \sum_{n\ge0}\frac{x^n}{n!} , \quad
   \cos x = 1 - \frac{x^2}{2!} + \frac{x^4}{4!} - \cdots , \quad
   \sin x = x - \frac{x^3}{3!} + \frac{x^5}{5!} - \cdots .
\end{equation*}
Putting $x = \mathrm{i}\theta$ in the first, the even powers
$(\mathrm{i}\theta)^{2m} = (-1)^m\theta^{2m}$ give the series of $\cos\theta$
and the odd powers $(\mathrm{i}\theta)^{2m+1} = \mathrm{i}(-1)^m
\theta^{2m+1}$ give $\mathrm{i}$ times that of $\sin\theta$:
\begin{equation*}
   \mathrm{e}^{\mathrm{i}\theta} = \cos\theta + \mathrm{i}\sin\theta ,
\end{equation*}
Euler's formula, an arrow of length 1 at the angle $\theta$. With
$-\theta$ in place of $\theta$ it gives $\mathrm{e}^{-\mathrm{i}\theta} =
\cos\theta - \mathrm{i}\sin\theta$, and adding and subtracting the two,
\begin{equation*}
   \cos\theta = \frac{\mathrm{e}^{\mathrm{i}\theta} +
   \mathrm{e}^{-\mathrm{i}\theta}}{2} , \qquad
   \sin\theta = \frac{\mathrm{e}^{\mathrm{i}\theta} -
   \mathrm{e}^{-\mathrm{i}\theta}}{2\mathrm{i}} ,
\end{equation*}
so any cosine or sine is a sum of two such exponentials. Exponentials
multiply by adding their powers, $\mathrm{e}^a\mathrm{e}^b = \mathrm{e}^{a
+ b}$, for complex numbers as for real ones: multiplying the two series
and collecting the terms of total power $n$ gives $\sum_{l}a^lb^{n -
l}/l!(n - l)! = (a + b)^n/n!$ by the binomial expansion, which uses only
the rules of adding and multiplying. So $\mathrm{e}^{\mathrm{i}A}
\mathrm{e}^{\mathrm{i}B} = \mathrm{e}^{\mathrm{i}(A + B)}$, which written
out is the pair of addition formulae of \cref{eq:mo-addition}: turning by
$A$ and then by $B$ turns by $A + B$. In this language the sums of
\cref{eq:ob-sq} are the real and imaginary parts of one sum, $\mathcal{C}
+ \mathrm{i}\mathcal{S} = \sum_j\mathrm{e}^{\mathrm{i}\vb{G}\cdot\vb{r}_j}$,
and $S(\vb{G})$ is its squared length over $N$.
\end{toolbox}

\subsection*{The Fourier transform}

A signal $x(t)$ that repeats with period $t_{\mathrm p}$ is a sum of
cosines and sines that repeat with it (\cref{sec:md-ewald}), which by the
toolbox are sums of $\mathrm{e}^{2\pi\mathrm{i}mt/t_{\mathrm p}}$ for whole
numbers $m$, positive and negative, each at the frequency $m/t_{\mathrm
p}$:
\begin{equation*}
   x(t) = \sum_mc_m\,\mathrm{e}^{2\pi\mathrm{i}mt/t_{\mathrm p}} , \qquad
   c_m = \frac{1}{t_{\mathrm p}}\int_{-t_{\mathrm p}/2}^{t_{\mathrm p}/2}
   x(t)\,\mathrm{e}^{-2\pi\mathrm{i}mt/t_{\mathrm p}}\,\dd t ,
\end{equation*}
the coefficients found, as in Chapter~10, because the average of
$\mathrm{e}^{2\pi\mathrm{i}(m - m')t/t_{\mathrm p}}$ over a period is 1 for
$m = m'$ and 0 otherwise. A signal that does not repeat is the limit of a
long period. The frequencies $\nu_m = m/t_{\mathrm p}$ then lie a small
distance $\Delta\nu = 1/t_{\mathrm p}$ apart, and with $\hat x(\nu_m) =
t_{\mathrm p}c_m$ the series reads $x(t) = \sum_m\hat x(\nu_m)\,
\mathrm{e}^{2\pi\mathrm{i}\nu_mt}\,\Delta\nu$, a sum of thin strips that
becomes an integral over $\nu$ as $\Delta\nu$ shrinks
(\cref{sec:mo-integrating}), while the integral for $\hat x(\nu_m)$ comes
to cover all times:
\begin{equation}
   \hat x(\nu) = \int_{-\infty}^{\infty}x(t)\,\mathrm{e}^{-2\pi\mathrm{i}
   \nu t}\,\dd t , \qquad
   x(t) = \int_{-\infty}^{\infty}\hat x(\nu)\,\mathrm{e}^{2\pi\mathrm{i}
   \nu t}\,\dd\nu ,
   \label{eq:ob-ft}
\end{equation}
the Fourier transform and its inverse. $\hat x(\nu)$ is a complex number
for each frequency, whose length says how much of that frequency the
signal holds and whose angle gives the wave's phase. A signal observed
for a time $t_{\mathrm f}$ cannot separate frequencies closer than about $1/t_{\mathrm f}$, since
two waves that differ by less than that drift apart by less than a whole
cycle over the record: a chord recorded for \SI{50}{\milli\second} shows
its notes only if they differ by more than \SI{20}{\hertz}.

\subsection*{The discrete Fourier transform}

A run records a signal at $N_{\mathrm s}$ equally spaced times $t_n =
n\Dt$, $n = 0, \dots, N_{\mathrm s} - 1$, and the integral becomes a sum,
\begin{equation}
   X_k = \sum_{n=0}^{N_{\mathrm s}-1}x_n\,
   \mathrm{e}^{-2\pi\mathrm{i}kn/N_{\mathrm s}} , \qquad
   x_n = \frac{1}{N_{\mathrm s}}\sum_{k=0}^{N_{\mathrm s}-1}X_k\,
   \mathrm{e}^{2\pi\mathrm{i}kn/N_{\mathrm s}} ,
   \label{eq:ob-dft}
\end{equation}
the discrete Fourier transform, $X_k$ belonging to the frequency $\nu_k =
k/N_{\mathrm s}\Dt$. The inverse holds because of a geometric sum.
Putting the first sum, with the index $n'$, into the second gives
$\frac{1}{N_{\mathrm s}}\sum_{n'}x_{n'}\sum_kw^k$ with $w =
\mathrm{e}^{2\pi\mathrm{i}(n - n')/N_{\mathrm s}}$, and $\sum_{k=0}^{N_{\mathrm
s}-1}w^k = (1 - w^{N_{\mathrm s}})/(1 - w)$, where $w^{N_{\mathrm s}} =
\mathrm{e}^{2\pi\mathrm{i}(n - n')} = 1$; so the sum over $k$ is 0 unless
$w = 1$, that is $n' = n$, when it is $N_{\mathrm s}$, and what remains is
$x_n$. For a real signal, changing $k$ to $N_{\mathrm s} - k$ only turns
each $\mathrm{e}^{-2\pi\mathrm{i}kn/N_{\mathrm s}}$ the other way, so
$X_{N_{\mathrm s}-k} = \bar X_k$: the upper half of the transform mirrors
the lower and belongs to the negative frequencies, and a cosine of height
1 puts $N_{\mathrm s}/2$ at each of its two places. Summed directly,
\cref{eq:ob-dft} needs $N_{\mathrm s}$ products for each of the
$N_{\mathrm s}$ values of $k$, $N_{\mathrm s}^2$ in all; the function
\texttt{dft} of \texttt{mdlab.analysis.spectra} does this.
\Cref{fig:ob-fourier}(a) shows $2\lvert X_k\rvert/400$, which is the
height of each wave, for the notes A (\SI{440}{\hertz}) and E
(\SI{659.25}{\hertz}) of height 1, sampled 8000 times a second for
\SI{50}{\milli\second}. The 400 samples are padded with zeros to 512, the
power of 2 that the fast transform below needs, whose grid of $8000/512 =
\SI{15.6}{\hertz}$ puts the two peaks at the frequencies nearest the
notes, 437.5 and \SI{656.25}{\hertz}.

\subsection*{The fast Fourier transform}

The sum splits into its even and odd samples, $n = 2p$ and $n = 2p + 1$:
\begin{equation*}
   X_k = \sum_{p=0}^{N_{\mathrm s}/2-1}x_{2p}\,\mathrm{e}^{-2\pi\mathrm{i}kp/(N_{\mathrm s}/2)}
   + \mathrm{e}^{-2\pi\mathrm{i}k/N_{\mathrm s}}\sum_{p=0}^{N_{\mathrm s}/2-1}x_{2p+1}\,
   \mathrm{e}^{-2\pi\mathrm{i}kp/(N_{\mathrm s}/2)}
   = X^{\mathrm e}_k + \mathrm{e}^{-2\pi\mathrm{i}k/N_{\mathrm s}}
   X^{\mathrm o}_k ,
\end{equation*}
with $X^{\mathrm e}$ and $X^{\mathrm o}$ the transforms of the even and the
odd samples, each of length $N_{\mathrm s}/2$. These repeat after $N_{\mathrm s}/2$ values of $k$,
and $\mathrm{e}^{-2\pi\mathrm{i}(k + N_{\mathrm s}/2)/N_{\mathrm s}} =
-\mathrm{e}^{-2\pi\mathrm{i}k/N_{\mathrm s}}$, so for $k < N_{\mathrm s}/2$
\begin{equation*}
   X_k = X^{\mathrm e}_k + \mathrm{e}^{-2\pi\mathrm{i}k/N_{\mathrm s}}
   X^{\mathrm o}_k , \qquad
   X_{k+N_{\mathrm s}/2} = X^{\mathrm e}_k - \mathrm{e}^{-2\pi\mathrm{i}k/N_{\mathrm s}}
   X^{\mathrm o}_k :
\end{equation*}
two transforms of half the length and $N_{\mathrm s}$ further operations
give the whole. When $N_{\mathrm s}$ is a power of 2 the halving repeats
down to transforms of one sample, each its own transform, and the work
$W(N_{\mathrm s}) = 2W(N_{\mathrm s}/2) + N_{\mathrm s}$ adds $N_{\mathrm s}$
at each of $\log_2N_{\mathrm s}$ levels, $N_{\mathrm s}\log_2N_{\mathrm s}$
in all. This is the fast
Fourier transform of Cooley and Tukey\cite{CooleyTukey1965}, the one
Chapter~10 promised for particle mesh Ewald. In \texttt{mdlab}:

\begin{code}
def fft(x: ArrayLike) -> NDArray:
    x = np.asarray(x, dtype=complex)
    n = len(x)
    if n == 1:
        return x
    if n % 2:
        raise ValueError("the length must be a power of 2")
    even, odd = fft(x[0::2]), fft(x[1::2])
    twiddle = np.exp(-2j * np.pi * np.arange(n // 2) / n) * odd
    return np.concatenate([even + twiddle, even - twiddle])
\end{code}

\noindent \texttt{dft} and \texttt{fft} agree with NumPy's FFT to
\num{4e-12} and \num{6e-14} on 256 samples. From $N_{\mathrm s} = 64$ to
4096, in three timings on the same machine, the direct sum takes about two
thousand times as long, half the 4096 by which $N_{\mathrm s}^2$ grows,
since a fixed cost of each call weighs most at small $N_{\mathrm s}$; the
recursive \texttt{fft} about 65 times, half the 128 of $N_{\mathrm
s}\log_2N_{\mathrm s}$, for the same reason; and NumPy's compiled FFT
about 10 times (\cref{fig:ob-fourier}c). Every later spectrum uses NumPy's.

\subsection*{Sampling and aliasing}

Samples taken every $\Dt$ cannot tell the frequency $\nu$ from $1/\Dt
- \nu$: at $t_n = n\Dt$,
\begin{equation*}
   \cos\bigl(2\pi(1/\Dt - \nu)n\Dt\bigr) = \cos(2\pi n - 2\pi\nu n
   \Dt) = \cos(2\pi\nu n\Dt) ,
\end{equation*}
since a whole number of turns changes no cosine (the sine changes sign).
The highest frequency that samples can show without confusion is half the
sampling rate, $\nu_{\mathrm N} = 1/2\Dt$, the Nyquist frequency; a
higher one appears folded back below it, at its alias. Films show this
when the spokes of a turning wheel seem to rotate slowly backwards. The E
of the chord, sampled 1000 times a second, gives samples that lie equally
on a wave of $1000 - 659.25 = \SI{340.75}{\hertz}$ (\cref{fig:ob-fourier}b),
to \num{1e-15}. For a run the consequence is direct: frames saved every
$\Dt$ cannot show vibrations faster than $1/2\Dt$, and a vibration
faster than that appears at a false frequency.

\subsection*{Correlation functions by the transform}

The products of \cref{eq:ob-dft} also compute correlations. For real
records $x$ and $y$, the conjugate of the transform of $x$ is $\bar X_k =
\sum_nx_n\mathrm{e}^{2\pi\mathrm{i}kn/N_{\mathrm s}}$, and
\begin{equation*}
   \frac{1}{N_{\mathrm s}}\sum_k\bar X_kY_k\,
   \mathrm{e}^{2\pi\mathrm{i}kj/N_{\mathrm s}}
   = \sum_n\sum_mx_ny_m\,\frac{1}{N_{\mathrm s}}\sum_k
   \mathrm{e}^{2\pi\mathrm{i}k(n + j - m)/N_{\mathrm s}}
   = \sum_nx_ny_{n+j} ,
\end{equation*}
since by the same geometric sum the inner sum is 1 when $m = n + j$ and 0
otherwise, with $n + j$ taken round the end of the record, as if the
record repeated. Padding each record with $N_{\mathrm s}$ zeros first makes
the wrapped products zero, leaving the sum over the $N_{\mathrm s} - j$
true pairs at the lag $j$; dividing by $N_{\mathrm s} - j$ gives the
average. Three transforms of length $2N_{\mathrm s}$ then do the work of
the $N_{\mathrm s}^2/2$ products of every pair. \texttt{stats.autocorrelation} of Chapter~15 and
\texttt{transport.correlation} of this chapter work this way.

\begin{figure}[htb]
\centering
\includegraphics{ch16_observables/fourier}
\caption{(a) The notes A (\SI{440}{\hertz}) and E (\SI{659.25}{\hertz})
of equal strength, sampled 8000 times a second for
\SI{50}{\milli\second} and padded to 512 samples: $2\lvert X_k\rvert/400$,
the height of each wave, against the frequency. (b) The E alone (line) sampled 1000 times a
second (points): the samples lie on a wave of \SI{340.75}{\hertz} as well
(dashed), its alias. (c) The time taken by the direct sum, by the
recursive radix-2 \texttt{fft} and by NumPy's FFT against the number of
samples, each relative to its own time at $N_{\mathrm s} = 64$, the least
of five tries.}
\label{fig:ob-fourier}
\end{figure}

\begin{keyidea}
The Fourier transform writes a signal as waves $\mathrm{e}^{2\pi\mathrm{i}
\nu t}$; a record of length $t_{\mathrm f}$ resolves frequencies about $1/t_{\mathrm f}$ apart.
The discrete transform of $N_{\mathrm s}$ samples costs $N_{\mathrm s}^2$
summed directly and $N_{\mathrm s}\log_2N_{\mathrm s}$ by halving, the fast
Fourier transform. Samples every $\Dt$
show frequencies only up to $1/2\Dt$, and fold higher ones back below
it. Correlation functions are computed through transforms of records
padded with zeros.
\end{keyidea}

\notebook{16}{sample a note and watch it fold back past the Nyquist frequency}

\begin{selfcheck}
A vibration of \SI{30}{\tera\hertz} is recorded with frames every
\SI{20}{\femto\second}. Where does it appear in the spectrum of the frames?
\end{selfcheck}
